\subsection{SUBGROUPS OF MAXIMAL RANK}

Recall (Chap. VI, §1, no. 7) that a subset P of R = R(G, T) is said to be closed if \((P + P) \cap R \subset P\) , and symmetric if P = -P.

PROPOSITION 12. Let H be the set of connected closed subgroups of G containing T, ordered by inclusion. The map \(\mathrm{H} \mapsto \mathrm{R}(\mathrm{H}, \mathrm{T})\) is an increasing bijection from H to the set of symmetric closed subsets of \(\mathrm{R}(\mathrm{G}, \mathrm{T})\) , ordered by inclusion.

If \(H \in H\) , then \(\mathrm{L}(\mathrm{H})_{(\mathbf{C})}\) is the direct sum of \(t_{C}\) and the \(g^{\alpha}\) for \(\alpha \in \mathrm{R}(\mathrm{H}, \mathrm{T})\) ; since this is a reductive subalgebra in \(g_{C}\) , the subset \(\mathrm{R}(\mathrm{H}, \mathrm{T})\) of R satisfies the stated conditions (Chap. VIII, §3, no. 1, Lemma 2 and Prop. 2). Conversely, if P is a subset of R satisfying these conditions, then \(t_{C} \oplus \sum_{\alpha \in P} g^{\alpha}\) is a subalgebra of \(g_{C}\) (loc. cit.) which is rational over R (no. 3), and hence of the form \(h_{(C)}\) , where h is a subalgebra of g. Let \(\mathrm{H}(\mathrm{P})\) be the integral subgroup of G defined by h; it is closed (§2, no. 4, Remark 1). We verify immediately that the maps \(\mathrm{H} \mapsto \mathrm{R}(\mathrm{H}, \mathrm{T})\) and \(\mathrm{P} \mapsto \mathrm{H}(\mathrm{P})\) are increasing and inverses of each other.

COROLLARY 1. There are only finitely-many closed subgroups of G containing T.

Let H be such a subgroup; then \(H_{0} \in H\) , and H is finite. Moreover, H is a subgroup of \(\mathrm{N}_{\mathrm{G}}(\mathrm{H}_{0})\) containing \(H_{0}\) , and \(\mathrm{N}_{\mathrm{G}}(\mathrm{H}_{0})/\mathrm{H}_{0}\) is finite ( \(\S2\) , no. 4, Prop. 4 and Remark 2).

COROLLARY 2. Let H be a connected closed subgroup of G containing T, and let \(W_{\mathrm{G}}^{\mathrm{H}}(\mathrm{T})\) be the stabilizer in \(W_{\mathrm{G}}(\mathrm{T})\) of the subset \(\mathrm{R}(\mathrm{H},\mathrm{T})\) of R. The group \(N_{\mathrm{G}}(\mathrm{H})/\mathrm{H}\) is isomorphic to the quotient group \(W_{\mathrm{G}}^{\mathrm{H}}(\mathrm{T})/W_{\mathrm{H}}(\mathrm{T})\) .

Indeed, it follows from Prop. 7 of §2, no. 5, applied to \(\mathrm{N}_{\mathrm{G}}(\mathrm{H})\) , that \(\mathrm{N}_{\mathrm{G}}(\mathrm{H})/\mathrm{H}\) is isomorphic to \(\mathrm{W}_{\mathrm{N}(\mathrm{H})}(\mathrm{T})/\mathrm{W}_{\mathrm{H}}(\mathrm{T})\) , where \(\mathrm{W}_{\mathrm{N}(\mathrm{H})}(\mathrm{T})\) is the set of elements of \(\mathrm{W}_{\mathrm{G}}(\mathrm{T})\) whose representatives in \(\mathrm{N}_{\mathrm{G}}(\mathrm{T})\) normalize H. Let \(n \in \mathrm{N}_{\mathrm{G}}(\mathrm{T})\) , and let w be its class in \(\mathrm{W}_{\mathrm{G}}(\mathrm{T})\) . By Chap. III, §9, no. 4, Prop. 11, n normalizes H if and only if \((\operatorname{Ad} n)(\operatorname{L}(\operatorname{H})) = \operatorname{L}(\operatorname{H})\) ; in view of Prop. 5 of no. 3, this also means that the subset \(\operatorname{R}(\operatorname{H}, \operatorname{T})\) of R is stable under w, hence the corollary.

Remark 1. The group \(W_{\mathrm{G}}^{\mathrm{H}}(\mathrm{T})\) is also the stabilizer in \(W_{\mathrm{G}}(\mathrm{T})\) of the subgroup \(\mathrm{C}(\mathrm{H})\) of T: this follows from Prop. 8 of no. 4.

PROPOSITION 13. Let H be a connected closed subgroup of G of maximal rank, and C its centre. Then C contains the centre of G, and H is the identity component of the centralizer of C.

Let S be a maximal torus of H. Since the centre of G is contained in S, it is contained in C. Put \(L = Z(C)_{0}\) ; this is a connected closed subgroup of G containing H, hence is of maximal rank, and its centre is equal to C. Denote by \(R_{H}\) and \(R_{L}\) the root systems of H and L, respectively, relative to S; then \(R_{H} \subset R_{L} \subset R(G, S)\) . Since \(C(H) = C(L)\) , Prop. 8 (no. 4) implies the equality \(Q(R_{H}) = Q(R_{L})\) ; but \(Q(R_{H}) \cap R_{L} = R_{H}\) (Chap. VI, §1, no. 7, Prop. 23), so \(R_{H} = R_{L}\) and H = L (Prop. 12).

Remark 2. Say that a subgroup \(\mathrm{C}\) of \(\mathrm{G}\) is radical if there exists a maximal torus \(\mathrm{S}\) of \(\mathrm{G}\) and a subset \(\mathrm{P}\) of \(\mathrm{R}(\mathrm{G},\mathrm{S})\) such that \(\mathrm{C} = \bigcap_{\alpha \in \mathrm{P}}\mathrm{Ker}\alpha\) . It follows

from Prop. 13 and Lemma 2 of no. 5 that the map \(H \mapsto C(H)\) induces a bijection from the set of connected closed subgroups of maximal rank to the set of radical subgroups of G. The inverse bijection is the map \(C \mapsto Z(C)_{0}\) .

COROLLARY. The set of \(g \in \mathbf{G}\) such that \(\mathrm{T} \cap g\mathrm{T}g^{-1} \neq \mathrm{C}(\mathrm{G})\) is the union of a finite number of closed analytic submanifolds of \(\mathbf{G}\) distinct from \(\mathbf{G}\) .

Indeed, put \(A_{g} = T \cap gTg^{-1}\) ; we have \(T \subset Z(A_{g})\) and \(gTg^{-1} \subset Z(A_{g})\) . Hence, there exists \(x \in Z(A_{g})\) such that \(xTx^{-1} = gTg^{-1}\) (§2, no. 2, Th. 2), which implies that \(g \in Z(A_{g}).N_{G}(T)\) . Denote by A the finite (Cor. 1) set of closed subgroups of G containing T and distinct from G, and put \(X = \bigcup_{H \in A} H.N_{G}(T)\) ; this is a finite union of closed submanifolds of G, distinct from G. If \(A_{g} \neq C(G)\) , then \(Z(A_{g}) \in A\) , and g belongs to X. Conversely, if \(g \in H.N_{G}(T)\) , with \(H \in A\) , then \(A_{g}\) contains \(C(H)\) , so \(A_{g} \neq C(G)\) (Prop. 13).

PROPOSITION 14. Let X be a subset of T, and let \(R_{X}\) be the set of roots \(\alpha \in \mathrm{R}(\mathrm{G}, \mathrm{T})\) such that \(\alpha(\mathrm{X}) = \{1\}\) . The group \(\mathrm{Z}_{\mathrm{G}}(\mathrm{X}) / \mathrm{Z}_{\mathrm{G}}(\mathrm{X})_{0}\) is isomorphic to the quotient of the subgroup of \(\mathrm{W}_{\mathrm{G}}(\mathrm{T})\) fixing X by the subgroup generated by the reflections \(s_{\alpha}\) for \(\alpha \in R_{X}\) .

Put \(H = Z_{\mathrm{G}}(\mathrm{X})\) ; since \(\mathrm{L}(\mathrm{H})_{(\mathbf{C})}\) is the set of points of \(g_{C}\) fixed by \(\mathrm{Ad}(\mathrm{X})\) , it is the sum of \(t_{C}\) and the \(g^{\alpha}\) for which \(\alpha(\mathrm{X}) = \{1\}\) . Consequently, \(\mathrm{R}(\mathrm{H}_{0}, \mathrm{T}) = \mathrm{R}_{\mathrm{X}}\) , so \(\mathrm{W}_{\mathrm{H}_{0}}(\mathrm{T})\) is generated by the reflections \(s_{\alpha}\) for \(\alpha \in R_{X}\) . It now suffices to apply Prop. 7 of §2, no. 5.

We shall see below ( \(\S5\) , no. 3, Th. 1) that if G is simply-connected and X reduces to a point, the centralizer \(Z(X)\) is connected.

